\documentclass[a4paper,12pt]{article}
\usepackage[polish]{babel}
\usepackage[T1]{fontenc}
\usepackage[utf8]{inputenc}
\usepackage{geometry}
\usepackage{enumitem}
\usepackage{hyperref}
\usepackage{listings}
\usepackage{xcolor}

\geometry{margin=2.5cm}

\definecolor{codebg}{RGB}{245,245,247}
\definecolor{codeframe}{RGB}{210,210,214}

\lstset{
  basicstyle=\ttfamily\small,
  backgroundcolor=\color{codebg},
  frame=single,
  rulecolor=\color{codeframe},
  breaklines=true,
  columns=fullflexible,
  showstringspaces=false,
  framexleftmargin=4pt,
  xleftmargin=4pt,
  literate=%
    {ą}{{\k a}}1 {ę}{{\k e}}1 {ć}{{\'c}}1 {ł}{{\l}}1 {ń}{{\'n}}1
    {ó}{{\'o}}1 {ś}{{\'s}}1 {ź}{{\'z}}1 {ż}{{\.z}}1
    {Ą}{{\k A}}1 {Ę}{{\k E}}1 {Ć}{{\'C}}1 {Ł}{{\L}}1 {Ń}{{\'N}}1
    {Ó}{{\'O}}1 {Ś}{{\'S}}1 {Ź}{{\'Z}}1 {Ż}{{\.Z}}1
}

\title{JSON Insight \\ \large Dokumentacja użytkownika}
\author{Filip Stępień \and Rafał Grot \and Bartłomiej Karkoszka}
\date{}

\begin{document}

\maketitle
\tableofcontents
\newpage

\section{Wprowadzenie}

\textbf{JSON Insight} to aplikacja okienkowa służąca do analizy, klasyfikacji oraz
przeszukiwania dokumentów w formacie JSON. Każdy dodany dokument jest sprawdzany pod
kątem poprawności składni, a następnie na podstawie jego struktury generowany jest
schemat (JSON Schema). Dokumenty o tej samej strukturze grupowane są w \emph{kolekcje},
co pozwala porządkować dane bez ręcznego ich opisywania. Zgromadzone dokumenty można
przeszukiwać przy użyciu prostego języka zapytań przypominającego SQL.

\section{Wymagania}

\begin{itemize}[noitemsep]
  \item \textbf{JDK 21} -- aplikacja kompiluje się i działa na Javie w wersji 21.
  \item \textbf{Maven} -- nie wymaga instalacji; w repozytorium znajduje się wrapper
        (\texttt{mvnw} / \texttt{mvnw.cmd}).
  \item System operacyjny: Linux, macOS lub Windows.
\end{itemize}

Dane aplikacji przechowywane są lokalnie w bazie SQLite w katalogu domowym użytkownika:

\begin{lstlisting}
~/.jsoninsight/jsoninsight.db
\end{lstlisting}

Katalog oraz tabele tworzone są automatycznie przy pierwszym uruchomieniu --
nie jest wymagana żadna ręczna konfiguracja bazy.

\section{Uruchomienie}

Z katalogu głównego projektu należy wykonać:

\begin{lstlisting}
# macOS / Linux
./mvnw clean javafx:run

# Windows
mvnw.cmd clean javafx:run
\end{lstlisting}

Polecenie pobierze zależności, skompiluje projekt i uruchomi okno aplikacji.
Klasą startową jest \texttt{io.github.jsoninsight.ui.JsonInsightApplication}.

\subsection{Pozostałe polecenia}

\begin{lstlisting}
./mvnw test            # testy jednostkowe (JUnit 5)
./mvnw clean compile   # sama kompilacja
\end{lstlisting}

\section{Układ interfejsu}

Główne okno składa się z czterech obszarów:

\begin{enumerate}[noitemsep]
  \item \textbf{Pasek narzędzi} (góra) -- przyciski dodawania i eksportu, pole zapytania,
        przycisk kreatora oraz przełącznik trybu ciemnego.
  \item \textbf{Panel boczny} (lewa strona) -- lista \emph{Kolekcji} oraz lista
        \emph{Dokumentów}.
  \item \textbf{Panel główny} (prawa strona) -- \emph{Podgląd dokumentu} (góra) oraz
        \emph{Wyniki wyszukiwania} ze stronicowaniem (dół).
  \item \textbf{Pasek statusu} (dół) -- komunikaty o stanie i wykonanych operacjach.
\end{enumerate}

\section{Funkcje}

\subsection{Dodawanie dokumentów}

Dokumenty można dodać na dwa sposoby:

\begin{itemize}[noitemsep]
  \item przyciskiem \textbf{,,Dodaj plik''}, który otwiera okno wyboru pliku \texttt{.json};
  \item metodą \textbf{przeciągnij i upuść} -- jeden lub wiele plików \texttt{.json}
        można upuścić bezpośrednio na okno aplikacji.
\end{itemize}

Po wczytaniu plik jest walidowany. Jeżeli zawartość nie jest poprawnym dokumentem JSON,
wyświetlany jest komunikat o błędzie, a plik nie zostaje zapisany. Dla poprawnego
dokumentu generowany jest schemat i porównywany z istniejącymi kolekcjami:

\begin{itemize}[noitemsep]
  \item przy zgodności struktury dokument trafia do \textbf{istniejącej kolekcji};
  \item w przeciwnym razie aplikacja prosi o nazwę dla \textbf{nowej kolekcji}
        (z proponowaną nazwą domyślną).
\end{itemize}

\subsection{Kolekcje}

Lista \emph{Kolekcje} prezentuje nazwę każdej kolekcji wraz z liczbą przypisanych do niej
dokumentów. Wybranie kolekcji:

\begin{itemize}[noitemsep]
  \item filtruje listę \emph{Dokumenty} do tej kolekcji,
  \item wyświetla jej schemat w panelu podglądu,
  \item zawęża zakres kolejnych wyszukiwań do tej kolekcji.
\end{itemize}

Nazwę kolekcji można zmienić, klikając ją prawym przyciskiem myszy i wybierając
\textbf{,,Zmień nazwę\ldots''}.

\subsection{Podgląd dokumentu}

Zaznaczenie dokumentu (na liście dokumentów lub w wynikach wyszukiwania) wyświetla jego
sformatowaną zawartość w panelu podglądu. Skrótem \texttt{Cmd/Ctrl + F} można wyszukać
i podświetlić fragment tekstu w podglądzie.

\subsection{Wyszukiwanie zapytaniem}

W polu zapytania na pasku narzędzi można wpisać:

\begin{itemize}[noitemsep]
  \item \textbf{sam predykat}, np. \texttt{.age >= 18 AND .active == true} --
        zostanie automatycznie owinięty w \texttt{SELECT * FROM \ldots WHERE \ldots};
  \item \textbf{pełne zapytanie}, np.
        \texttt{SELECT .name, .age FROM uzytkownicy WHERE .age > 21}.
\end{itemize}

Zakres przeszukiwania zależy od zaznaczenia: przy wybranej kolekcji przeszukiwane są
tylko jej dokumenty, w przeciwnym razie -- wszystkie. Wyniki są stronicowane
(20 pozycji na stronę), a nawigację zapewniają przyciski poprzedniej i następnej
strony. Pełny opis składni języka zapytań znajduje się
w pliku \texttt{doc/query-language.md}.

\subsection{Kreator zapytań}

Przycisk \textbf{,,Kreator''} otwiera wizualny budowniczy zapytań, który nie wymaga
znajomości składni. Pozwala on:

\begin{itemize}[noitemsep]
  \item dodawać pojedyncze \textbf{warunki} (pole, operator, wartość) oraz zagnieżdżone
        \textbf{grupy} warunków,
  \item łączyć elementy operatorami \textbf{AND} / \textbf{OR} i negować całe grupy
        przyciskiem \textbf{NOT},
  \item korzystać z operatorów: porównań (\texttt{==}, \texttt{!=}, \texttt{>},
        \texttt{>=}, \texttt{<}, \texttt{<=}), \texttt{EXISTS} / \texttt{NOT EXISTS},
        \texttt{IS STRING/NUMBER/BOOLEAN/NULL/ARRAY/OBJECT}, dopasowania wyrażeniem
        regularnym (\texttt{matches}) oraz rozmiaru (\texttt{size}),
  \item korzystać z podpowiedzi nazw pól zbieranych z istniejących dokumentów,
  \item obserwować podgląd zapytania na żywo, sprawdzić jego poprawność przyciskiem
        \textbf{,,Sprawdź''} i przekazać je do wyszukiwarki przyciskiem
        \textbf{,,Zastosuj''}.
\end{itemize}

\subsection{Eksport danych}

\begin{itemize}[noitemsep]
  \item \textbf{Pobierz dokument} -- zapisuje zaznaczony dokument do pliku \texttt{.json}.
  \item \textbf{Pobierz schemat} -- zapisuje (sformatowany) schemat zaznaczonej kolekcji
        do pliku \texttt{.json}.
  \item \textbf{Eksport wyników} -- zapisuje wszystkie wyniki wyszukiwania do pliku
        \texttt{.jsonl} lub \texttt{.csv}; format wybierany jest rozszerzeniem pliku.
\end{itemize}

\subsection{Tryb ciemny}

Przełącznik \textbf{,,Dark''} na pasku narzędzi zmienia motyw kolorystyczny aplikacji.
Ustawienie obejmuje również okno kreatora zapytań.

\section{Typowy scenariusz pracy}

\begin{enumerate}[noitemsep]
  \item Uruchom aplikację poleceniem \texttt{./mvnw clean javafx:run}.
  \item Dodaj pliki JSON przyciskiem ,,Dodaj plik'' lub przeciągając je na okno;
        nadaj nazwy nowym kolekcjom.
  \item Przeglądaj kolekcje i dokumenty, korzystając z panelu bocznego oraz podglądu.
  \item Wyszukaj interesujące dokumenty -- wpisując zapytanie ręcznie lub budując je
        w kreatorze.
  \item Wyeksportuj pojedynczy dokument, schemat kolekcji lub komplet wyników
        wyszukiwania.
\end{enumerate}

\section{Struktura projektu}

\begin{lstlisting}
src/main/java/io/github/jsoninsight/
  ui/        # JavaFX: aplikacja, kontrolery, widoki FXML
  service/   # uslugi: dokumenty, schematy, zapytania (+ impl/)
  model/     # model danych: JsonDocument, Collection, JsonSchema...
  json/      # lekser/parser JSON, generator schematu
  query/     # jezyk zapytan: lekser, parser, AST, ewaluator
\end{lstlisting}

\paragraph{Technologie.} JavaFX 21, SQLite (sqlite-jdbc), Gson, Lombok, JUnit~5, Maven.

\end{document}
